\section{Challenges}
\label{sec:challenges}

Designing a programmable policy substrate that spans the GPU driver and device execution layer introduces unique challenges not present in CPU-only eBPF environments. These challenges stem from the monolithic nature of GPU drivers, the massive parallelism of the SIMT execution model, and the disjoint memory spaces between host and device.

\subsection{C1: Safe and Efficient Driver Extensibility}
Modern GPU drivers (e.g., NVIDIA's open kernel modules) are massive, complex codebases that tightly couple resource management policies with low-level hardware interactions. Unlike the Linux kernel, which has evolved stable tracepoints and \texttt{struct\_ops} interfaces for eBPF, GPU drivers lack designed-for-programmability hooks.
\begin{itemize}[leftmargin=*]
    \item \textbf{Mechanism-Policy Coupling:} Decisions about page migration, prefetching, and context scheduling are often buried deep within hardware-specific state machines. Retrofitting programmability requires identifying "narrow waist" points where policy can be decoupled from mechanism without destabilizing the driver or incurring excessive overhead on critical paths (e.g., interrupt handlers).
    \item \textbf{Safety Risks:} Naive injection of logic into these paths is dangerous. A buggy policy that holds a lock too long or corrupts a hardware command queue can cause the entire GPU to hang, requiring a node reboot. The challenge is to expose meaningful control without exposing the raw hardware interface to potentially unsafe user code.
\end{itemize}

\subsection{C2: Safe and Efficient Device-Side Execution}
Extending eBPF-style programmability to the GPU device itself is fundamentally difficult because the execution model is radically different from the CPU.
\begin{itemize}[leftmargin=*]
    \item \textbf{SIMT Mismatch:} Standard eBPF assumes a scalar, single-threaded execution model. GPUs execute threads in warps (groups of 32 threads) that execute in lockstep. Naively running a scalar eBPF program on every thread would cause massive divergence and serialization, destroying performance.
    \item \textbf{Memory Safety in a Parallel World:} The GPU memory hierarchy (global, shared, local) and consistency models are complex. A policy running on the device must not only be memory-safe in isolation but also thread-safe when accessing shared state concurrently with thousands of other warps. The standard eBPF verifier does not model SIMT execution or GPU memory consistency.
    \item \textbf{No OS on Device:} There is no kernel on the GPU to handle exceptions or enforce limits. If a policy enters an infinite loop or accesses invalid memory, it can hang the streaming multiprocessor (SM). We need a way to enforce safety guarantees (bounded execution, valid memory access) statically, before code is ever loaded onto the device.
\end{itemize}

\subsection{C3: Unified Cross-Layer State and Management}
Effective policies often require coordination between the host driver and the device. For example, a prefetching policy needs to know which pages are being accessed on the device (Device signal) to trigger migration on the host (Host action).
\begin{itemize}[leftmargin=*]
    \item \textbf{Disjoint Address Spaces:} Host and device have separate physical memory and virtual address spaces. Sharing state (e.g., a BPF map) typically requires explicit, slow PCIe transfers or complex pointer swizzling.
    \item \textbf{Asynchronous Execution:} The host driver runs asynchronously from the GPU. Correlating events (e.g., "which kernel caused this page fault?") is difficult because the GPU pipeline is deep and opaque. A unified substrate must provide a coherent view of time and state across this asynchronous boundary without stalling the GPU.
\end{itemize}
